\section{Camino euleriano}
\whenuse{Recorrer cada arista exactamente una vez; tambi\'en cuando objetos son
aristas entre dos tipos de entidad, como filas y columnas.}

Hierholzer usa cada arista exactamente una vez. El chequeo final de longitud
tamb\'ien valida que todas las aristas pertenezcan al componente recorrido.

\complexity{$O(V+E)$ tiempo y memoria.}

\begin{lstlisting}[style=cpp]
vector<int> directed_eulerian_path(
        const vector<vector<pair<int,int>>> &graph, int edge_count) {
    int n = graph.size();
    if (n == 0) return edge_count == 0 ? vector<int>{} : vector<int>{};
    vector<int> indegree(n, 0), outdegree(n, 0);
    for (int v = 0; v < n; ++v)
        for (auto [to, id] : graph[v]) {
            (void)id;
            ++outdegree[v]; ++indegree[to];
        }

    int start = -1, finish = -1;
    for (int v = 0; v < n; ++v) {
        int difference = outdegree[v] - indegree[v];
        if (difference == 1) {
            if (start != -1) return {};
            start = v;
        } else if (difference == -1) {
            if (finish != -1) return {};
            finish = v;
        } else if (difference != 0) return {};
    }
    if ((start == -1) != (finish == -1)) return {};
    if (start == -1) {
        start = 0;
        while (start < n && outdegree[start] == 0) ++start;
        if (start == n) return {0};
    }

    vector<int> next_edge(n, 0), stack{start}, path;
    vector<char> used(edge_count, false);
    while (!stack.empty()) {
        int v = stack.back();
        while (next_edge[v] < (int)graph[v].size() &&
               used[graph[v][next_edge[v]].second]) ++next_edge[v];
        if (next_edge[v] == (int)graph[v].size()) {
            path.push_back(v); stack.pop_back();
        } else {
            auto [to, id] = graph[v][next_edge[v]++];
            if (id < 0 || id >= edge_count) return {};
            used[id] = true;
            stack.push_back(to);
        }
    }
    if ((int)path.size() != edge_count + 1) return {};
    reverse(path.begin(), path.end());
    return path;
}

vector<int> undirected_eulerian_path(
        const vector<vector<pair<int,int>>> &graph, int edge_count) {
    int n = graph.size();
    if (n == 0) return {};
    vector<int> odd;
    for (int v = 0; v < n; ++v)
        if (graph[v].size() % 2 == 1) odd.push_back(v);
    if (!odd.empty() && odd.size() != 2) return {};

    int start = odd.empty() ? 0 : odd[0];
    while (start < n && graph[start].empty()) ++start;
    if (start == n) return {0};

    vector<int> next_edge(n, 0), stack{start}, path;
    vector<char> used(edge_count, false);
    while (!stack.empty()) {
        int v = stack.back();
        while (next_edge[v] < (int)graph[v].size() &&
               used[graph[v][next_edge[v]].second]) ++next_edge[v];
        if (next_edge[v] == (int)graph[v].size()) {
            path.push_back(v); stack.pop_back();
        } else {
            auto [to, id] = graph[v][next_edge[v]++];
            if (id < 0 || id >= edge_count) return {};
            used[id] = true;
            stack.push_back(to);
        }
    }
    if ((int)path.size() != edge_count + 1) return {};
    reverse(path.begin(), path.end());
    return path;
}
\end{lstlisting}

En el grafo no dirigido, cada arista debe aparecer dos veces en la lista de
adyacencia con el mismo ID. En el dirigido aparece una sola vez.
